\documentclass[11pt,a4paper]{article}

\usepackage[margin=2.4cm]{geometry}
\usepackage{amsmath,amssymb}
\usepackage{booktabs}
\usepackage{xcolor}
\usepackage{tcolorbox}
\usepackage{fancyhdr}
\usepackage{enumitem}
\usepackage[hidelinks]{hyperref}

\tcbuselibrary{skins,breakable}

\definecolor{keyblue}{HTML}{1F4E79}
\definecolor{warnamber}{HTML}{8A6100}
\definecolor{softgrey}{HTML}{F4F5F7}
\definecolor{rulegrey}{HTML}{C9CDD3}

\pagestyle{fancy}
\fancyhf{}
\lhead{\small CS F407 --- Search and A*}
\rhead{\small Personal study notes}
\cfoot{\small\thepage}
\renewcommand{\headrulewidth}{0.4pt}

\newtcolorbox{keyidea}[1]{
  breakable, colback=softgrey, colframe=keyblue,
  boxrule=0.9pt, left=8pt, right=8pt, top=6pt, bottom=6pt,
  title={\bfseries #1}, fonttitle=\small\bfseries
}

\newtcolorbox{warnbox}[1]{
  breakable, colback=softgrey, colframe=warnamber,
  boxrule=0.9pt, left=8pt, right=8pt, top=6pt, bottom=6pt,
  title={\bfseries #1}, fonttitle=\small\bfseries
}

\setlist[itemize]{itemsep=2pt, topsep=4pt}
\setlist[enumerate]{itemsep=2pt, topsep=4pt}

\title{\vspace{-1.5cm}\textbf{Search and A*: An End-to-End Explainer}}
\author{Warehouse Robot Navigation --- CS F407 Artificial Intelligence}
\date{Study notes, not for submission}

\begin{document}
\maketitle

\begin{keyidea}{How to read this}
This is the whole laboratory explained from first principles: what a search
problem is, how BFS and A* differ, what actually happens inside the code, and why
the experiments came out the way they did. Every number quoted here was measured
from the programs in this folder --- nothing is illustrative.

If you only read one section, read \S6. It is the part of this laboratory that
teaches something you cannot get from a textbook.
\end{keyidea}

\tableofcontents
\vspace{6pt}
\hrule

%==========================================================================
\section{The problem}

A robot sits in a warehouse laid out as a grid. It starts at \texttt{S}, must
reach \texttt{G}, and shelves (\texttt{\#}) block the way. It moves one cell at a
time --- up, down, left or right --- and every move costs 1.

\begin{verbatim}
        #################
        #S....#.........#          S = start  (row 1, column 1)
        #.###.#.#######.#          G = goal   (row 7, column 15)
        #...#.#.......#.#          # = shelf
        ###.#.#######.#.#          . = free floor
        #...#.........#.#
        #.###########.#.#          9 rows x 17 columns
        #.............#G#          64 free cells
        #################
\end{verbatim}

The answer, which the program finds, is a 40-move route. But the point of the
laboratory is not the number 40. It is \emph{how} you get a machine to find it,
and how you convince yourself the machine is right.

%==========================================================================
\section{Turning a maze into a search problem}

Search is a general method. Before you can apply it, you have to say what the
problem \emph{is} in the vocabulary the method uses. That vocabulary has six
parts:
\[
  \mathcal{P} = (S,\; A,\; T,\; s_0,\; G,\; c)
\]

\begin{center}
\begin{tabular}{@{}lll@{}}
\toprule
\textbf{Symbol} & \textbf{Name} & \textbf{For this warehouse} \\
\midrule
$S$   & states           & every free cell, as a pair \texttt{(row, col)} --- 64 of them \\
$A$   & actions          & Up, Down, Left, Right \\
$T$   & transition       & move one cell, if the target is in bounds and not a shelf \\
$s_0$ & initial state    & \texttt{(1, 1)} \\
$G$   & goal states      & \texttt{\{(7, 15)\}} \\
$c$   & cost             & 1 per move \\
\bottomrule
\end{tabular}
\end{center}

Most of this is obvious. One part is not, and it is the part that matters.

\subsection{State versus node --- the idea everything depends on}

A \textbf{state} is a configuration of the world. A \textbf{node} is a state
\emph{plus} the bookkeeping the search is carrying about how it got there.

For this problem the state is just the cell, \texttt{(row, col)}. It is tempting
to also put ``the route I took to get here'' or ``how many steps I have used''
into the state. Do not. Those are \emph{history}, and they belong to the node.

Why it matters: suppose two different routes both arrive at cell \texttt{(5,7)},
one having used 12 moves and one having used 18. From \texttt{(5,7)} onward, the
remaining problem is \emph{identical} in both cases --- the same walls, the same
goal, the same choices. That is what lets the search say ``I have already been
here, and I got here more cheaply last time, so I can discard this.''

If you folded the route into the state, those two arrivals would be two
\emph{different} states, nothing would ever look like a repeat, and the search
would degenerate into enumerating every possible walk through the maze --- which
is astronomically many.

\begin{keyidea}{The test for what belongs in a state}
Ask: \emph{if I told you only this, could you work out what the robot can do next
and what it will cost?} If yes, it is state. If it only tells you about the past,
it is history --- keep it on the node.
\end{keyidea}

So in the code:
\begin{itemize}
  \item \texttt{state} is \texttt{(r, c)} --- a plain tuple.
  \item \texttt{g}, the parent pointer, and \texttt{f} live in dictionaries keyed
        by state. They are node information.
\end{itemize}

\subsection{Why a tuple}

\texttt{(r, c)} is a tuple, not a list, for two concrete reasons: tuples are
\textbf{hashable}, so they can be dictionary keys and set members --- which is
exactly what ``have I seen this state before?'' needs; and they are
\textbf{immutable}, so a state cannot be accidentally changed after being stored
as a key.

%==========================================================================
\section{The family of search algorithms}

Every algorithm in this family shares one loop:

\begin{enumerate}
  \item Take a state out of the \textbf{frontier} (the set of places you know
        about but have not explored from yet).
  \item If it is the goal, stop.
  \item Otherwise \textbf{expand} it: generate its legal neighbours and put them
        in the frontier.
\end{enumerate}

\begin{keyidea}{The only thing that differs between these algorithms}
They differ in \emph{one} decision: which state to take out of the frontier next.
That single choice is what separates BFS from DFS from A*. Everything else is the
same loop.
\end{keyidea}

\begin{center}
\begin{tabular}{@{}llll@{}}
\toprule
\textbf{Algorithm} & \textbf{Picks next} & \textbf{Frontier} & \textbf{Optimal?} \\
\midrule
BFS            & shallowest (fewest moves)      & FIFO queue     & yes, if all costs equal \\
DFS            & deepest                        & LIFO stack     & no \\
Uniform-cost   & cheapest so far, $g$           & priority queue & yes \\
Greedy         & looks closest to goal, $h$     & priority queue & no \\
A*             & best total estimate, $g + h$   & priority queue & yes, if $h$ admissible \\
\bottomrule
\end{tabular}
\end{center}

\subsection{Blind versus informed}

BFS, DFS and uniform-cost are \textbf{blind} (or uninformed): they know where
they have been, but nothing about where the goal is. BFS expands in a growing
ring around the start, treating the direction of the goal as irrelevant.

A* is \textbf{informed}: it is given a function $h$ that estimates the cost still
remaining. That is the entire difference, and the word ``informed'' means exactly
this and nothing more.

%==========================================================================
\section{\texorpdfstring{$f = g + h$}{f = g + h}}

A* orders its frontier by
\[
  f(n) = g(n) + h(n)
\]
where
\begin{itemize}
  \item $g(n)$ = cost actually paid to reach $n$ from the start --- \emph{known},
  \item $h(n)$ = estimated cost from $n$ to the goal --- \emph{guessed},
  \item $f(n)$ = estimated total cost of a solution passing through $n$.
\end{itemize}

So A* always works on the node that currently looks like it lies on the best
complete solution. Note the shape of this: $g$ looks backwards and is a fact;
$h$ looks forwards and is an estimate. A* is the algorithm that adds a fact to a
guess and sorts by the sum.

\subsection{The heuristic used here}

\[
  h\big((r,c)\big) = |r - r_G| + |c - c_G| \qquad \text{(Manhattan distance)}
\]

Named after the street grid: you cannot cut diagonally through a block, so the
distance is the sum of the horizontal and vertical legs.

Where does it come from? \textbf{Relax the problem.} Take the real warehouse and
delete every shelf. In that easier world, the cheapest route from $(r,c)$ to the
goal is exactly $|r-r_G| + |c-c_G|$ moves --- you close the row gap and the
column gap, and no move closes both at once because diagonals are not allowed.

\begin{keyidea}{Where good heuristics come from}
Solve an easier version of your problem exactly, and use that answer as the
estimate. Because you removed constraints, the easier answer can never be worse
than the real one --- which is precisely the property A* needs.
\end{keyidea}

\subsection{Worked example on the real map}

Start is $(1,1)$, goal is $(7,15)$.
\[
  h(s_0) = |1-7| + |1-15| = 6 + 14 = 20
\]
So \textbf{no route can be shorter than 20 moves}. The program returns 40 --- the
shelves double the ideal distance. That 20 is a genuinely useful check: if the
program had ever returned 18, it would be wrong without needing to look at the path.

Take a cell partway along, say $(5,9)$, reached in $g = 12$ moves:
\[
  h = |5-7| + |9-15| = 2 + 6 = 8, \qquad f = 12 + 8 = 20.
\]

%==========================================================================
\section{What the code actually does}

\subsection{The three data structures}

\begin{center}
\begin{tabular}{@{}lll@{}}
\toprule
\textbf{Name} & \textbf{Type} & \textbf{Holds} \\
\midrule
\texttt{frontier} & min-heap  & states to explore, ordered by $f$ \\
\texttt{best\_g}  & dict      & cheapest known cost to each state \\
\texttt{parent}   & dict      & which state we came from \\
\texttt{closed}   & set       & states already expanded \\
\bottomrule
\end{tabular}
\end{center}

\subsection{Why a heap}

You need ``give me the smallest $f$'' repeatedly. Scanning a list is $O(n)$ every
time. A binary heap gives you the minimum in $O(\log n)$. Python's
\texttt{heapq} compares tuples left to right, so pushing $(f, \ldots, \text{state})$
sorts by $f$ automatically.

\subsection{Why there is a counter in the tuple}

Entries are pushed as \texttt{(f, -g, counter, state)}. The counter exists because
if $f$ and $-g$ both tie, Python moves on to the next element and would start
comparing the \emph{state tuples} --- silently ordering the search by geography,
which is meaningless. A unique, increasing counter guarantees the comparison
always stops there.

\subsection{Lazy deletion --- and why ``generated'' exceeds ``expanded''}

If we find a cheaper route to a state already sitting in the frontier, we ought to
lower its priority. Python's heap has no such operation. So instead we push a
\emph{second} entry for that state with the better $f$, and let both sit there.

The better one is nearer the top, so it comes out first. When the stale one
eventually surfaces, the state is already in \texttt{closed} and we skip it.

\begin{keyidea}{Reading the statistics}
\textbf{Generated} counts pushes onto the frontier. \textbf{Expanded} counts
states actually taken out and explored. Generated is always $\geq$ expanded,
because some pushes are stale duplicates that get skipped. This is not a bug ---
it is what lazy deletion looks like from the outside.
\end{keyidea}

\subsection{Testing for the goal on \emph{pop}, not on generation}

This is the subtlest correctness point in the whole program.

When you generate a neighbour and notice it is the goal, it is tempting to return
immediately. \textbf{Do not.} At that moment you have found \emph{a} route to the
goal, not necessarily the \emph{cheapest} one --- a better route may still be
sitting in the frontier, unexplored.

A* is only guaranteed optimal if you wait until the goal is \emph{popped}. When
the goal comes off the top of the heap, its $f$ is the smallest in the frontier,
and since $h(\text{goal}) = 0$, that $f$ is its true cost $g$. Nothing cheaper
remains. Only then is it safe to stop.

%==========================================================================
\section{The bug that taught the lesson}

This is the important section.

\subsection{What happened}

The generated A* was correct. It passed all 20 tests. It returned optimal paths
on every map. And then this measurement appeared, on an open test map with 130
free cells:

\begin{center}
\begin{tabular}{@{}lrrr@{}}
\toprule
& \textbf{Path length} & \textbf{States expanded} & \textbf{Peak frontier} \\
\midrule
BFS               & 25 & 116 & 8 \\
A* with Manhattan & 25 & 116 & 8 \\
\bottomrule
\end{tabular}
\end{center}

Identical. Not similar --- \emph{identical}, including the peak frontier size.
An informed search does not coincidentally match a blind one to the digit. The
heuristic was being computed correctly and then having no effect whatsoever.

\subsection{Why}

Consider any cell on a sensible, detour-free route from start to goal. Every move
along such a route increases $g$ by 1 and decreases $h$ by 1. So the sum does not
move:
\[
  f = g + h = \text{constant} = h(s_0).
\]

Checking the open map directly: \textbf{all 130 free cells have $f = 25$.} There
is no ordering information in $f$ at all. It is one enormous plateau.

\begin{warnbox}{The failure}
When every entry ties on $f$, the heap's ordering is decided entirely by the
tie-break. The tie-break was first-in-first-out. FIFO order over a frontier
\emph{is} breadth-first order. So A* was running BFS --- correctly, optimally,
and pointlessly.
\end{warnbox}

\subsection{The fix}

Break ties toward the \textbf{larger $g$} --- the node further along a route. Push
$(f, -g, \text{counter}, \text{state})$; the heap wants small values, and $-g$ is
smallest when $g$ is largest.

Intuition: among many nodes that all claim the same total cost 25, prefer the one
that has already \emph{paid} 22 of it over one that has paid 3. The first is
nearly at the goal; the second has all the work still ahead. Ties should be
resolved toward finishing, not toward fanning out.

\subsection{The result, drawn from real traces}

Below, \texttt{*} marks a cell the algorithm actually expanded and \texttt{.}
marks one it never touched. These are recorded traces, not sketches.

\vspace{4pt}
\noindent\textbf{BFS --- 116 of 130 cells (89\% of the map)}
\begin{verbatim}
        #####################
        #S******************#
        #*******************#
        #****#########******#
        #****#.......#******#
        #****#.......#******#
        #****#########******#
        #*******************#
        #******************G#
        #####################
\end{verbatim}

\noindent\textbf{A* with the tie-break fixed --- 26 of 130 cells (20\%)}
\begin{verbatim}
        #####################
        #S..................#
        #*..................#
        #*...#########......#
        #*...#.......#......#
        #*...#.......#......#
        #*...#########......#
        #*..................#
        #*****************.G#
        #####################
\end{verbatim}

BFS floods the room. A* walks down the left wall and along the bottom --- almost
exactly the path, and nearly nothing else. Same 25-move answer, 78\% less work.

\subsection{The lesson}

\begin{keyidea}{Working output $\neq$ validated algorithm}
No test caught this. Every path returned was legal and optimal. The defect was
found because a \emph{number was implausible} --- and it was only implausible
because there was an expectation to compare it against.

If you accept generated code because its output looks right, you inherit exactly
this class of defect: silently inert machinery, invisible precisely where the
code is most sophisticated.
\end{keyidea}

Note also where the fault really lay. The design specified the frontier
completely except for the tie-break rule, which it left vague. The LLM
implemented what was specified and chose the obvious default for what was not.
\textbf{It was a good engineer working to bad drawings.}

%==========================================================================
\section{Why the supplied warehouse hides all of this}

On the warehouse itself, BFS and A* both expand all 64 cells. A* saves nothing.
This is not a failure --- it is a property of the map.

The peak frontier on the warehouse is \textbf{3}. That means at almost every step
the search had at most three options. The warehouse is a maze of single-width
corridors: you walk until you hit a junction, and there are very few junctions.

\begin{keyidea}{When a heuristic can help}
A heuristic \emph{ranks alternatives}. If there are no alternatives, there is
nothing to rank. In a corridor, informed and blind search do exactly the same
thing, because there is only one thing to do.
\end{keyidea}

This is why the report added an open-plan control map. Without it, the honest
conclusion would have been ``A* gave no benefit,'' which is true of this map and
badly misleading about A*.

%==========================================================================
\section{Three properties of a heuristic}

These get confused constantly. They are different things.

\subsection{Admissible: never overestimates}
\[
  h(n) \leq h^*(n) \quad\text{for every } n
\]
where $h^*$ is the true remaining cost. Manhattan is admissible here: it is the
answer to the shelf-free problem, and putting shelves back can only make the real
route longer.

\textbf{Admissibility buys the optimality guarantee.} If $h$ is admissible, A*
returns a shortest path.

\subsection{Consistent: never jumps}
\[
  h(n) \leq c(n, n') + h(n')
\]
Each move changes Manhattan distance by at most 1, and each move costs 1 --- so
it holds. Consistency is stronger than admissibility.

\textbf{Consistency is what makes the closed set safe.} With a consistent $h$,
the first time you expand a state you already have its cheapest route, so you
never need to reopen it.

\subsection{Informative: how close to the truth}

An admissible heuristic can still be useless. $h = 0$ is perfectly admissible and
tells you nothing. The closer $h$ gets to $h^*$ from below, the more sharply it
discriminates and the fewer states you expand.

\begin{center}
\begin{tabular}{@{}lrrl@{}}
\toprule
\textbf{Heuristic (open map)} & \textbf{Length} & \textbf{Expanded} & \\
\midrule
$h = 0$          & 25 & 116 & admissible, useless \\
Euclidean        & 25 &  92 & admissible, weak \\
Manhattan        & 25 &  26 & admissible, sharp \\
\bottomrule
\end{tabular}
\end{center}

All three optimal. A factor of 4.5 between them in work done.

Why is Euclidean weaker? $\sqrt{dr^2 + dc^2} \leq |dr| + |dc|$ always. Euclidean
is the true distance only if you could fly diagonally --- and you cannot, so it
underestimates by more than it needs to. It is a looser lower bound, and looser
lower bounds discriminate less.

\begin{keyidea}{Admissible and informative are independent}
Admissibility is about \emph{correctness}: never overestimate. Informativeness is
about \emph{speed}: get as close to the truth as you can. $h=0$ has the first and
none of the second.
\end{keyidea}

%==========================================================================
\section{Breaking admissibility on purpose}

What if $h$ \emph{does} overestimate? Try $h = 2 \times$ Manhattan.

On the warehouse and the open map, it returned optimal paths anyway. That is
worth understanding: inadmissibility \textbf{removes the guarantee; it does not
force a violation}. So the experiment was pushed further --- a randomised search
over maps for a case where it actually breaks. On a cluttered room:

\begin{center}
\begin{tabular}{@{}lrr@{}}
\toprule
\textbf{Heuristic} & \textbf{Path length} & \textbf{States expanded} \\
\midrule
Manhattan          & 28 & 133 \\
$2 \times$ Manhattan & \textbf{30} & \textbf{31} \\
\bottomrule
\end{tabular}
\end{center}

77\% less work --- and a path two moves worse than optimal. Repeated over 2{,}754
random solvable maps, doubling the heuristic returned a suboptimal path on
\textbf{25.9\%} of them, worst case 8 moves too long.

\subsection{Why overestimating breaks it}

A* stops when the goal is popped, reasoning that anything still in the frontier
has $f \geq f(\text{goal})$ and therefore cannot be better. That reasoning depends
on $f$ being a \emph{lower bound}. If $h$ overestimates, a genuinely cheaper route
can be sitting in the frontier carrying an inflated $f$ that makes it look worse
than it is. A* never gets to it, and stops with the worse answer.

\begin{keyidea}{The trade-off, stated plainly}
Inflating the heuristic makes A* more goal-directed and much faster, and pays for
it with the optimality guarantee. That is sometimes a good trade --- in a game
with a millisecond budget, a path 7\% too long delivered 4$\times$ faster is a
clear win. The point is that it must be a \emph{decision}, not an accident.
\end{keyidea}

%==========================================================================
\section{What to remember}

\begin{enumerate}
  \item \textbf{State versus node.} State is the cell. History belongs to the
        node. Getting this wrong breaks duplicate detection and the search
        explodes.

  \item \textbf{$f = g + h$.} A known cost plus an estimated one. $g$ backwards
        and factual, $h$ forwards and guessed.

  \item \textbf{Test the goal on pop, not on generation.} Otherwise you return
        the first route found, not the best one.

  \item \textbf{Admissible $\Rightarrow$ optimal.} Never overestimate and A* is
        guaranteed to find a shortest path.

  \item \textbf{Admissible does not mean useful.} $h=0$ is admissible and
        worthless. Closeness to $h^*$ is what buys speed: 116 / 92 / 26.

  \item \textbf{A* only helps where there are choices.} In a corridor it degrades
        to BFS, not because it is broken but because there is nothing to rank.

  \item \textbf{Ties matter as much as the heuristic.} On a unit-cost grid
        Manhattan makes $f$ constant, and the tie-break rule alone decides whether
        A* expands 26 states or 116.

  \item \textbf{Working output is not a validated algorithm.} The bug here
        produced correct, optimal answers on every test while doing nothing. It
        was caught by an expectation, not by a test.
\end{enumerate}

%==========================================================================
\section{Questions worth being able to answer}

\begin{description}[style=unboxed,leftmargin=0pt,itemsep=6pt]

\item[Why is position alone enough to be the state?]
Because the warehouse is static and the robot carries nothing. Position fully
determines what moves are available and what they cost. Everything else is
history.

\item[Why test the goal on pop rather than on generation?]
Generating the goal only proves \emph{a} route exists. When the goal is popped,
its $f$ is minimal in the frontier and $h=0$, so $f = g$ is its true cost and
nothing cheaper remains.

\item[Why is Manhattan the right heuristic here?]
It is the exact solution to the relaxed problem with no shelves. Adding shelves
back can only lengthen routes, so it never overestimates --- admissible, and also
consistent since each unit move changes it by at most 1.

\item[Why did A* expand as many states as BFS on the warehouse?]
Peak frontier 3 --- the map is single-width corridor with almost no junctions.
There were no alternatives for the heuristic to rank.

\item[Why can generated exceed expanded?]
Lazy deletion. An improved route to a state is pushed as a second heap entry
rather than updating the existing one, and the stale entry is skipped on pop.

\item[What happens with $h = 0$?]
A* becomes uniform-cost search, which with unit costs behaves like BFS. Still
optimal, maximum work.

\item[What happens if $h$ overestimates?]
The optimality guarantee is lost. A cheaper route can carry an inflated $f$, be
overlooked, and A* stops on a worse path --- measured here at 25.9\% of random
maps, with a 77\% saving in expansions when it happens.

\item[If BFS is also optimal here, why bother with A*?]
Because BFS is optimal only when every step costs the same, and it does far more
work: 116 expansions against 26 on the open map. Change the costs --- carpet that
slows the robot, say --- and BFS stops being optimal at all, while A* with an
admissible heuristic still is.

\end{description}

\vspace{10pt}
\hrule
\vspace{6pt}
\noindent\small\textit{All figures measured from} \texttt{astar\_warehouse.py},
\texttt{test\_search.py} \textit{and} \texttt{experiments.py}\textit{. Verified
against an independently written BFS oracle over 6{,}000 random maps ---
36{,}000 A* runs, zero disagreements.}

\end{document}
